% Generated draft: requires mathematical and source review.
\documentclass[11pt]{article}
\usepackage[margin=25mm]{geometry}
\usepackage{fontspec}
\setmainfont{Latin Modern Roman}
\usepackage{amsmath,amssymb,mathtools,unicode-math}
\setmathfont{Latin Modern Math}
\usepackage{enumitem,needspace,xurl,hyperref}
\setlength{\parindent}{0pt}
\setlength{\parskip}{5pt}
\setlength{\emergencystretch}{3em}
\newcommand{\N}{\mathbb{N}}
\newcommand{\Z}{\mathbb{Z}}
\newcommand{\Q}{\mathbb{Q}}
\newcommand{\R}{\mathbb{R}}
\newcommand{\C}{\mathbb{C}}
\newcommand{\F}{\mathbb{F}}
\newcommand{\A}{\mathbb{A}}
\newcommand{\PP}{\mathbb{P}}
\newcommand{\E}{\mathbb{E}}
\newcommand{\eps}{\varepsilon}
\newcommand{\dd}{\,\mathrm{d}}
\begin{document}
\section*{UnsolvedMath ID: 2000009}
\textbf{Source problem code:} catalogue entry
\subsection{Definitions and mathematical statement}
The problem is to devise general existence methods for local but nonglobal minimizers and other weak equilibria in nonlinear elastostatics.

To separate the relevant notions, let $\mathcal A$ be a specified class of admissible deformations, let $I:\mathcal A\to\R\cup\{+\infty\}$ be an elastic total-energy functional, and fix a topology $\tau$ on $\mathcal A$. A deformation $y\in\mathcal A$ with finite energy is a $\tau$-local minimizer if some $\tau$-neighborhood $U$ of $y$ satisfies $I(y)\leq I(v)$ for all $v\in U\cap\mathcal A$. It is nonglobal if
\[
 \exists v\in\mathcal A,\qquad I(v)<I(y).
\]
The choice of topology, for example a strong Sobolev or uniform topology, changes the local-minimum requirement and is not specified in the source.

A standard contextual energy has the form
\[
 I(y)=\int_\Omega W(x,\nabla y(x))\,\dd x-\mathcal L(y),
\]
with stored-energy density $W$, loading functional $\mathcal L$, and admissible boundary/geometric conditions built into $\mathcal A$. When differentiability, integrability and admissible variations justify it, a weak equilibrium satisfies
\[
 \int_\Omega D_FW(x,\nabla y):\nabla\varphi\,\dd x
 =D\mathcal L(y)[\varphi]
\]
for all admissible test variations $\varphi$. This formula is conditional context, not an assertion that all singular or constrained elastic models admit this first-variation description.

The requested method should establish existence of such nonglobal local minima or other weak equilibria under meaningful hypotheses in nonlinear elasticity. The source supplies no fixed dimension, domain, stored-energy class, growth/singularity assumptions, admissible-space definition, boundary/loading data, topology of local minimality, or precise generality/completeness criterion. Hence no single universal existence theorem is determined, and none of those hypotheses is silently invented.
\subsection{Short English statement}
Find methods that produce physically admissible local energy minima that fail to be global minima, and further weak equilibria, in nonlinear elastic boundary-value problems. The supplied moderate-confidence OPEN-TRIAGE assessment mentions isolated multiple-equilibrium and bifurcation constructions, including necking, shear-band, symmetry-breaking and multiple-well examples, but identifies a general method as unestablished. It refers to broad polyconvex/quasiconvex energy classes without giving precise hypotheses. Those examples and the status claim have no specific supplied supporting bibliography and are not independently verified here.
\subsection{Sources}
\begin{itemize}
\item UnsolvedMath permanent ID 2000009, source code catalogue entry, ``Some Open Problems in Elasticity --- Nonglobal local minimizers.'' Registry curation and metadata: Ulam AI and UnsolvedMath Contributors, CC BY 4.0; updated 2026-08-06.
\item Ball, source list ``problems on Elasticity'' (2002), Problem 9: \url{https://people.maths.ox.ac.uk/ball/Articles%20in%20Conference%20Proceedings%20and%20Books/JMB%202002%20re%20Marsden%2060th.pdf}. Supplied PDF-text extraction accessed 2026-07-29. The metadata says the source presents an open item; status/release and rights review are NEEDS\_REVIEW.
\item The supplied moderate-confidence research assessment provides no specific references for its cited example mechanisms or its general-status conclusion. No missing author, title, or theorem is inferred.
\end{itemize}
\end{document}
